\scanpage{27}
\uline{Prop 6} \deleted{Supposons $A$ \unsure{non-atomique}} \inserted{Soit $A\in\underline b^+$, on a}
\begingroup\csname tagsleft@true\endcsname
\[
\begin{aligned}
\varphi_A^-(t)&=\mathop{\rm Inf}\limits_{\substack{E\in P(\mathcal A)\\\mathrm{Tr}\,E<t}}
\mathop{\rm Sup}\limits_{x\in(1-E).H}(Ax,x)\\
&=\mathop{\rm Inf}\limits_{\substack{E\in P(\mathcal A)\\\mathrm{Tr}\,E<t}}\lVert(1-E)A(1-E)\rVert
\end{aligned}\tag*{(5)}
\]
\endgroup

\uline{Démonstration} Soit $\underline C$ \inserted{$=\mathcal C_0(M)$} une sous-algèbre commutative autoadjointe \unsure{maximale} dans $\underline b$, contenant $A$. Si on se borne dans \unsure{le calcul des|le calcul du} $\mathop{\rm Inf}$ aux $E$ qui sont dans $\underline C$, on trouve
\[
\mathop{\rm Inf}\limits_{\substack{E\ \text{\unsure{mesurable|ensemble}}\subset M\\\mathrm{Tr}(E)<t}}
\lVert(1-E)\deleted{\unsure{E}}A\rVert_\infty\unsure{\geq|=}\varphi_A^-(t)
\]
\figure{\begin{tikzpicture}[x=0.46cm,y=0.46cm,line cap=round]
\draw (0,3.2) .. controls (.32,2.42) and (.63,1.72) .. (.9,1.64);
\draw (.67,2.15) -- (1.0,1.72);
\draw (1.1,1.63) .. controls (1.65,1.57) and (2.25,1.47) .. (2.85,1.47);
\draw (2.13,1.51) -- (2.1,.06); \fill (2.1,.06) circle (.045);
\draw (3,1.13) -- (3.56,.53);
\end{tikzpicture}}
(N\textsuperscript{o} 1, \unsure{formule} (4)) \deleted{\unsure{d'où|d'ici}} \unsure{donc} dans (5) $\varphi^-_A(t)$ est $\geq$ le deuxième mb.

% slanted
\margin{on pose $\lvert A\rvert=\sqrt{A^*A}$} \uline{Corollaire 1} $A\in\underline b$ implique
\begingroup\csname tagsleft@true\endcsname
\[
\begin{aligned}
\varphi^-_{\lvert A\rvert}(t)&=\mathop{\rm Inf}\limits_{\substack{E\in P(\mathcal A)\\\unsure{\mathrm{Tr}|\mathrm T}(E)<t}}
\mathop{\rm Sup}\limits_{x\in(1-E)H}\lVert Ax\rVert\\
&=\mathop{\rm Inf}\limits_{\substack{E\in P(\mathcal A)\\T(\unsure{E|t})<t}}\lVert A(1-E)\rVert
\end{aligned}\tag*{(6)}
\]
\endgroup

\uline{Corollaire 2} $A\in\underline b$, $B\unsure{\in|\epsilon}\mathcal A$ impliquent
\begingroup\csname tagsleft@true\endcsname
\[
\begin{aligned}
&\left\{\begin{aligned}
\varphi_{\lvert BA\rvert}&\leq\lVert B\rVert\varphi_{\lvert A\rvert}\\
\varphi_{\lvert AB\rvert}&\unsure{\leq|=}\lVert B\rVert\varphi_{\lvert A\rvert}
\end{aligned}\right.
\end{aligned}\tag*{(7)}
\]
\endgroup
% slanted
\margin{\uline{Corollaire} 5 $\varphi_{\lvert1+A\rvert}\leq\varphi_{1+\lvert A\rvert}$}

\uline{Corollaire 3} $\varphi_A$ est fonction croissante de $A\in\underline b^+$

\uline{Corollaire 4} Les fonctions $A\unsure{\text{--}|\mapsto}\varphi^-_{\lvert A\rvert}(t)$ sur $b$ forment un ensemble équicontinu de fonctions (quand $\underline b$ est \unsure{muni} de sa topologie \unsure{normée|usuelle})
% archive sheet number: 28

